%! Linear Algebra — Cumulative Final Exam — Practice Final --- Study Copy (blank)
\documentclass[10pt]{article}
\usepackage{linalg-article}
\usepackage{linalg-boxes}
\newcommand{\vv}{\mathbf{v}}
\newcommand{\ww}{\mathbf{w}}
\newcommand{\xx}{\mathbf{x}}
\newcommand{\bb}{\mathbf{b}}
\newcommand{\pp}{\mathbf{p}}
\newcommand{\ee}{\mathbf{e}}
\newcommand{\zero}{\mathbf{0}}
\newcommand{\uu}{\mathbf{u}}
\newcommand{\avec}{\mathbf{a}}
\newcommand{\relu}{\mathrm{ReLU}}
\newcommand{\T}{\mathsf{T}}

% Part-divider strip
\newcommand{\parthead}[1]{%
  \vspace{6pt}\noindent
  \begin{headlinebox}{burgundy}{\color{white}\bfseries #1}\end{headlinebox}%
  \vspace{4pt}}

\begin{document}

\pageheader{Linear Algebra --- Cumulative Final}{Practice Final --- Study Copy}
\namedateperiod

\vspace{0.10in}
\begin{remindbox}
\small
\textbf{This is a practice final.} It covers the whole course (Units 1--8) and mirrors the
real final in length ($50$ questions), format, and the ideas it tests --- but the numbers and
contexts differ. Work every problem with no notes, then check your answers against the key.
\textbf{Parts:} A Vocabulary ($16$) $\cdot$ B Multiple Choice ($12$) $\cdot$ C Short Answer \&
Computation ($16$) $\cdot$ D Extended Response ($6$). \hfill\textbf{Total: 100 pts.}
\end{remindbox}

% ======================================================================
\parthead{Part A --- Vocabulary (16 pts)}
Write the letter of the best description next to each term. Two matching sets.

\vspace{2pt}
\noindent\textit{Set 1 --- Units 1--4.}\\[2pt]
\noindent
\begin{minipage}[t]{0.40\textwidth}
\begin{enumerate}[leftmargin=2.2em, itemsep=5pt, label=\arabic*.]
  \item Unit vector \blank{1.1cm}
  \item Dot product \blank{1.1cm}
  \item Pivot \blank{1.1cm}
  \item Inverse matrix \blank{1.1cm}
  \item Nullspace $N(A)$ \blank{1.1cm}
  \item Basis \blank{1.1cm}
  \item Projection \blank{1.1cm}
  \item Least squares \blank{1.1cm}
\end{enumerate}
\end{minipage}%
\hfill
\begin{minipage}[t]{0.57\textwidth}
\small
\begin{description}[leftmargin=1.4em, itemsep=2pt]
  \item[(A)] $A^{-1}$, the matrix that undoes $A$: $A^{-1}A=I$.
  \item[(B)] The set of all solutions to $A\xx=\zero$.
  \item[(C)] A vector whose length is exactly $1$.
  \item[(D)] An independent set that spans a space --- the fewest vectors that build every vector in it.
  \item[(E)] $\vv\cdot\ww=v_1w_1+v_2w_2$; it measures length and angle.
  \item[(F)] The first nonzero entry in a row, used to clear the entries below it.
  \item[(G)] Sending $\bb$ to the closest point $\pp$ in a subspace.
  \item[(H)] Choosing $\hat\xx$ to make the error $\|\bb-A\xx\|^2$ as small as possible.
\end{description}
\end{minipage}

\vspace{6pt}
\noindent\textit{Set 2 --- Units 5--8.}\\[2pt]
\noindent
\begin{minipage}[t]{0.40\textwidth}
\begin{enumerate}[leftmargin=2.2em, itemsep=5pt, label=\arabic*., start=9]
  \item Determinant \blank{1.1cm}
  \item Linear transformation \blank{1.1cm}
  \item Eigenvector \blank{1.1cm}
  \item Diagonalization \blank{1.1cm}
  \item Singular value \blank{1.1cm}
  \item Principal component \blank{1.1cm}
  \item ReLU \blank{1.1cm}
  \item Covariance matrix \blank{1.1cm}
\end{enumerate}
\end{minipage}%
\hfill
\begin{minipage}[t]{0.57\textwidth}
\small
\begin{description}[leftmargin=1.4em, itemsep=2pt]
  \item[(A)] The map $\xx\mapsto A\xx$ that fixes the origin and sends grid lines to straight, evenly spaced lines.
  \item[(B)] A number $\sigma\ge0$: how far $A$ stretches along a singular direction ($\sigma=\sqrt\lambda$ of $A^\T A$).
  \item[(C)] A nonzero vector $\xx$ whose direction is unchanged by $A$: $A\xx=\lambda\xx$.
  \item[(D)] The rectifier $\max(0,x)$: keep positive inputs, send negatives to $0$.
  \item[(E)] A number whose absolute value is the area (or volume) factor; it is $0$ exactly when the matrix is singular.
  \item[(F)] Writing $A=X\Lambda X^{-1}$ with the eigenvectors in $X$ and the eigenvalues down $\Lambda$.
  \item[(G)] The symmetric matrix of variances and covariances; its eigenvectors are the principal components.
  \item[(H)] A top singular direction of centered data --- the direction of greatest spread.
\end{description}
\end{minipage}

% ======================================================================
\parthead{Part B --- Multiple Choice (24 pts)}
Circle the best answer. ($2$ pts each.)

\begin{enumerate}[leftmargin=1.9em, itemsep=8pt]
  \item The length of the vector $(6,8)$ is
    \hfill (A) $10$ \quad (B) $14$ \quad (C) $48$ \quad (D) $100$
  \item Two nonzero vectors are perpendicular exactly when
    \hfill (A) they are parallel \quad (B) they have equal length \quad (C) their dot product is $0$ \quad (D) their sum is $\zero$
  \item If elimination reaches a zero in a pivot position with no row below to swap in, the matrix is
    \hfill (A) invertible \quad (B) singular \quad (C) symmetric \quad (D) orthogonal
  \item For an $m\times n$ matrix of rank $r$, the dimension of the nullspace $N(A)$ is
    \hfill (A) $r$ \quad (B) $m-r$ \quad (C) $n-r$ \quad (D) $n$
  \item The column space and the row space of a matrix always have
    \hfill (A) the same dimension $r$ \quad (B) dimensions $r$ and $n$ \quad (C) dimension $0$ \quad (D) no relation
  \item In a least-squares fit, the error vector $\ee=\bb-\pp$ is
    \hfill (A) in the column space \quad (B) perpendicular to the column space \quad (C) always $\zero$ \quad (D) parallel to $\bb$
  \item A square matrix has $\det=0$ exactly when it is
    \hfill (A) symmetric \quad (B) orthogonal \quad (C) singular (no inverse) \quad (D) diagonal
  \item For square matrices, $\det(AB)$ equals
    \hfill (A) $\det A+\det B$ \quad (B) $\det A\,\det B$ \quad (C) $\det A-\det B$ \quad (D) $1$
  \item The number $\lambda$ is an eigenvalue of $A$ when $A\xx=\lambda\xx$ holds for
    \hfill (A) $\xx=\zero$ \quad (B) some nonzero $\xx$ \quad (C) every $\xx$ \quad (D) no $\xx$
  \item A symmetric matrix $S=S^\T$ is guaranteed to have
    \hfill (A) complex eigenvalues \quad (B) no eigenvectors \quad (C) real eigenvalues and perpendicular eigenvectors \quad (D) equal eigenvalues
  \item Which is true of \emph{every} matrix $A$ (any shape)?
    \hfill (A) it is invertible \quad (B) it has an SVD $A=U\Sigma V^\T$ \quad (C) it is symmetric \quad (D) it has real eigenvalues
  \item Gradient descent updates the weights by stepping
    \hfill (A) along the slope (uphill) \quad (B) opposite the slope (downhill) \quad (C) to $\zero$ \quad (D) to a random point
\end{enumerate}

% ======================================================================
\parthead{Part C --- Short Answer \& Computation (48 pts)}
Show your work. ($3$ pts each.)

\begin{enumerate}[leftmargin=1.9em, itemsep=10pt]
  \item \textbf{(U1)} Let $\vv=(3,4)$ and $\ww=(4,-3)$.
        \textbf{(a)} Compute $2\vv+\ww$. \textbf{(b)} Compute $\vv\cdot\ww$ and the lengths $\|\vv\|,\|\ww\|$.
        Are $\vv$ and $\ww$ perpendicular?
        \vspace{2.2cm}
  \item \textbf{(U1)} Let $A=\begin{bsmallmatrix}1&2\\2&4\end{bsmallmatrix}$.
        \textbf{(a)} Compute $A\begin{bsmallmatrix}1\\3\end{bsmallmatrix}$.
        \textbf{(b)} Are the columns independent? What is the rank, and what does the column space look like?
        \vspace{2.2cm}
  \item \textbf{(U2)} Solve by elimination: $\;x+2y=5,\quad 3x+4y=11.$
        \vspace{2.2cm}
  \item \textbf{(U2)} Let $A=\begin{bsmallmatrix}2&1\\5&3\end{bsmallmatrix}$.
        \textbf{(a)} Find $\det A$. \textbf{(b)} Find $A^{-1}$ and check that $A^{-1}A=I$.
        \vspace{2.4cm}
  \item \textbf{(U3)} Reduce $A=\begin{bsmallmatrix}1&2&1\\2&4&3\end{bsmallmatrix}$ and find a special solution of $A\xx=\zero$.
        What is the dimension of the nullspace?
        \vspace{2.4cm}
  \item \textbf{(U3)} The $3\times4$ matrix $A=\begin{bsmallmatrix}1&2&0&1\\0&0&1&3\\0&0&0&0\end{bsmallmatrix}$.
        Give the rank $r$ and the dimensions of all four subspaces: $C(A)$, $N(A)$, the row space $C(A^\T)$, and $N(A^\T)$.
        \vspace{2.6cm}
  \item \textbf{(U4)} Project $\bb=(4,3)$ onto the line through $\avec=(1,2)$. \textbf{(a)} Find the projection $\pp$.
        \textbf{(b)} Find the error $\ee=\bb-\pp$ and verify it is perpendicular to $\avec$.
        \vspace{2.4cm}
  \item \textbf{(U4)} Find the best-fit line $b=C+Dt$ through the points $(0,6),(1,0),(2,0)$ using the normal
        equations $A^\T A\,\hat\xx=A^\T\bb$.
        \vspace{2.8cm}
  \item \textbf{(U5)} Compute the determinants.
        \textbf{(a)} $\det\begin{bsmallmatrix}3&1\\2&4\end{bsmallmatrix}$.
        \textbf{(b)} $\det\begin{bsmallmatrix}1&2&0\\0&3&1\\2&0&1\end{bsmallmatrix}$ (expand along a row).
        \vspace{2.4cm}
  \item \textbf{(U5)} The transformation $A=\begin{bsmallmatrix}2&1\\0&3\end{bsmallmatrix}$.
        \textbf{(a)} What area does a unit square map to? \textbf{(b)} Does $A$ preserve or flip orientation?
        \textbf{(c)} If $\det B=\tfrac12$, what does $AB$ do to area?
        \vspace{2.4cm}
  \item \textbf{(U6)} Let $A=\begin{bsmallmatrix}2&1\\1&2\end{bsmallmatrix}$. Find its eigenvalues and an eigenvector
        for each. Check your eigenvalues against the trace and determinant.
        \vspace{2.6cm}
  \item \textbf{(U6)} Let $A=\begin{bsmallmatrix}4&1\\2&3\end{bsmallmatrix}$. Use the trace and determinant to find the
        eigenvalues, find an eigenvector for each, and write the diagonalization $A=X\Lambda X^{-1}$ (give $X$ and $\Lambda$).
        \vspace{2.8cm}
  \item \textbf{(U7)} Let $A=\begin{bsmallmatrix}2&2\\1&-2\end{bsmallmatrix}$. Form $A^\T A$, find its eigenvalues,
        and give the singular values of $A$.
        \vspace{2.6cm}
  \item \textbf{(U7)} \textbf{(a)} Compute the outer product $\uu\vv^\T$ for $\uu=(1,2)$, $\vv=(3,1)$, and say why it is
        rank one. \textbf{(b)} The matrix $A=\begin{bsmallmatrix}5&3\\3&5\end{bsmallmatrix}$ has singular values $8,2$.
        What fraction of the energy does the top rank-one layer keep?
        \vspace{2.4cm}
  \item \textbf{(U8)} Evaluate one neural-network layer $\relu(A\vv+\bb)$ with
        $A=\begin{bsmallmatrix}1&1\\1&-1\end{bsmallmatrix}$, $\vv=\begin{bsmallmatrix}1\\3\end{bsmallmatrix}$,
        $\bb=\begin{bsmallmatrix}0\\1\end{bsmallmatrix}$. Show $A\vv$, then $A\vv+\bb$, then the ReLU.
        \vspace{2.4cm}
  \item \textbf{(U8)} Four students' two quiz scores are $(2,4),(4,2),(6,8),(8,6)$.
        \textbf{(a)} Find the mean point and the covariance matrix $V$.
        \textbf{(b)} Find the eigenvalues of $V$ and the fraction of the total variance the first principal
        component ($\mathrm{PC}_1$) holds.
        \vspace{2.8cm}
\end{enumerate}

% ======================================================================
\parthead{Part D --- Extended Response (12 pts)}
Explain your reasoning in full sentences. ($2$ pts each.)

\begin{enumerate}[leftmargin=1.9em, itemsep=12pt]
  \item \textbf{(Units 1--3 --- rank)} Explain what the rank $r$ of an $m\times n$ matrix counts (pivots
        $=$ independent columns), and how it fixes the dimensions of the four subspaces. Include why
        $\dim C(A)+\dim N(A)=n$.
        \vspace{2.4cm}
  \item \textbf{(Unit 4 --- least squares)} Explain why the best-fit solution $\hat\xx$ comes from
        \emph{projecting} $\bb$ onto the column space, and why the error must be perpendicular to $C(A)$ ---
        the idea behind the normal equations $A^\T A\,\hat\xx=A^\T\bb$.
        \vspace{2.4cm}
  \item \textbf{(Units 5--6 --- singular)} Explain the chain: $\det A=0$ $\iff$ $A$ is singular $\iff$ the
        columns are dependent $\iff$ $0$ is an eigenvalue. What does $\det A$ measure geometrically?
        \vspace{2.4cm}
  \item \textbf{(Unit 6 --- spectral theorem)} State the spectral theorem for a symmetric matrix $S=S^\T$.
        Why are its eigenvalues real and its eigenvectors perpendicular, and why does this let you write
        $S=Q\Lambda Q^\T$ with an orthogonal $Q$?
        \vspace{2.4cm}
  \item \textbf{(Unit 7 --- the SVD)} Explain why \emph{every} matrix $A$ has a singular value decomposition,
        building the argument from the fact that $A^\T A$ is symmetric (real $\lambda\ge0$, perpendicular
        eigenvectors) so $\sigma=\sqrt\lambda$. Why is this stronger than diagonalization by eigenvectors?
        \vspace{2.4cm}
  \item \textbf{(Unit 8 --- the whole course)} A neural network that learns from data is built entirely from
        tools in this course. Name the linear-algebra idea behind each stage: (i) a layer, (ii) scoring
        predictions with a loss, (iii) choosing the weights, (iv) understanding the shape of the data.
        \vspace{2.6cm}
\end{enumerate}

\end{document}
